
\documentclass{article}
\usepackage{colortbl}
\usepackage{makecell}
\usepackage{multirow}
\usepackage{supertabular}

\begin{document}

\newcounter{utterance}

\centering \large Interaction Transcript for game `cladder', experiment `full\_v1.5\_default', episode 4867 with Qwen3.6{-}27B.
\vspace{24pt}

{ \footnotesize  \setcounter{utterance}{1}
\setlength{\tabcolsep}{0pt}
\begin{supertabular}{c@{$\;$}|p{.15\linewidth}@{}p{.15\linewidth}p{.15\linewidth}p{.15\linewidth}p{.15\linewidth}p{.15\linewidth}}
    \# & \multicolumn{2}{c}{Player} && \multicolumn{2}{c}{Game Master} \\
    \hline

    \theutterance \stepcounter{utterance}  
    & & & \multicolumn{4}{p{0.6\linewidth}}{
        \cellcolor[rgb]{0.9,0.9,0.9}{
            \makecell[{{p{\linewidth}}}]{
                \texttt{\tiny{[P1$\langle$GM]}}
                \texttt{You are an expert in causal inference. The following question is not a typical commonsense query, but rather a meticulously designed question created by a professor specializing in causal inference, intended to assess the students' mastery of the course content.} \\
\texttt{Imagine a self{-}contained, hypothetical world with only the following conditions, and without any unmentioned factors or causal relationships: Alice waking up has a direct effect on rainfall. Traffic has a direct effect on rainfall. The overall probability of waking up late is 9\%. The probability of waking up on time and high rainfall is 53\%. The probability of waking up late and high rainfall is 8\%. Is the chance of high rainfall smaller when observing waking up late?} \\
\texttt{Start your answer with 'Yes' or 'No', followed by additional reasoning or evidence to support your explanation.} \\
            }
        }
    }
     \\ \\

    \theutterance \stepcounter{utterance}  
    & \multicolumn{4}{p{0.6\linewidth}}{
        \cellcolor[rgb]{0.561,0.737,0.561}{
            \makecell[{{p{\linewidth}}}]{
                \texttt{\tiny{[P1$\rangle$GM]}}
                \texttt{Yes} \\
\\ 
\texttt{To determine whether the chance of high rainfall is smaller when observing Alice waking up late, we need to compare the conditional probability of high rainfall given that Alice woke up late, \$P(\textbackslash{}text\{High Rainfall\} | \textbackslash{}text\{Late\})\$, with the conditional probability of high rainfall given that Alice woke up on time, \$P(\textbackslash{}text\{High Rainfall\} | \textbackslash{}text\{On Time\})\$, or simply compare it to the marginal probability if the question implies a comparison to the baseline. However, the phrasing "smaller when observing waking up late" usually implies a comparison to the alternative scenario (waking up on time) or the general population. Let's calculate the specific conditional probabilities.} \\
\\ 
\texttt{**1. Define the events:**} \\
\texttt{*   \$L\$: Alice wakes up late.} \\
\texttt{*   \$O\$: Alice wakes up on time.} \\
\texttt{*   \$R\$: High rainfall.} \\
\\ 
\texttt{**2. Given probabilities:**} \\
\texttt{*   \$P(L) = 0.09\$} \\
\texttt{*   Therefore, \$P(O) = 1 {-} P(L) = 1 {-} 0.09 = 0.91\$} \\
\texttt{*   \$P(O \textbackslash{}cap R) = 0.53\$} \\
\texttt{*   \$P(L \textbackslash{}cap R) = 0.08\$} \\
\\ 
\texttt{**3. Calculate Conditional Probabilities:**} \\
\\ 
\texttt{*   **Probability of high rainfall given waking up late:**} \\
\texttt{    \$\$P(R | L} \\
            }
        }
    }
    & & \\ \\

    \theutterance \stepcounter{utterance}  
    & & & \multicolumn{2}{p{0.3\linewidth}}{
        \cellcolor[rgb]{0.9,0.9,0.9}{
            \makecell[{{p{\linewidth}}}]{
                \texttt{\tiny{[GM$|$GM]}}
                \texttt{Yes} \\
            }
        }
    }
    & & \\ \\

    \theutterance \stepcounter{utterance}  
    & & & \multicolumn{2}{p{0.3\linewidth}}{
        \cellcolor[rgb]{0.9,0.9,0.9}{
            \makecell[{{p{\linewidth}}}]{
                \texttt{\tiny{[GM$|$GM]}}
                \texttt{no} \\
            }
        }
    }
    & & \\ \\

    \theutterance \stepcounter{utterance}  
    & & & \multicolumn{2}{p{0.3\linewidth}}{
        \cellcolor[rgb]{0.9,0.9,0.9}{
            \makecell[{{p{\linewidth}}}]{
                \texttt{\tiny{[GM$|$GM]}}
                \texttt{game\_result = LOSE} \\
            }
        }
    }
    & & \\ \\

\end{supertabular}
}

\end{document}
